\section{练习}

\begin{question}\begin{drill}[铁块从浮板移入水中]\label{s5:drill:1}
浮板托着质量$m$、密度$\rho_i>\rho$的铁块漂浮。将铁块移入容器并沉底，浮板仍漂浮，无溢流。容器截面积恒为$S$，求液面高度变化。
\end{drill}
\fig{fig:s5:drill:1:s}{题设：铁块从浮板移入水中}{\begin{tikzpicture}[scale=.88]\draw[world](0,3)--(0,0)--(4,0)--(4,3);\fill[blue!8](0,.05)rectangle(4,2);\draw[guide](0,2)--(4,2);\draw[axis](4.6,2.5)--(4.6,-.3)node[below]{$z$};\draw[world](1,1.8)rectangle(3,2.15);\draw[world](1.7,2.15)rectangle(2.3,2.8);\node[above]at(2,2.8){$m,\rho_i$};\node[left]at(1,1.85){浮板};\node at(.5,.7){$\rho$};\node at(2,-.5){容器截面积$S$};\end{tikzpicture}}
\end{question}
\answerlink{s5:drill:1}
\begin{question}\begin{drill}[竖直闸门的压力中心]\label{s5:drill:2}
宽$b$、高$h$的竖直矩形闸门，上边缘在深度$z_0$，一侧为密度$\rho$的水，另一侧大气。求合力及作用点深度。
\end{drill}
\fig{fig:s5:drill:2:s}{正视图定义门宽和门高，侧视图定义液面深度与湿润范围}{\begin{tikzpicture}

\begin{scope}
\draw[world](0,0)rectangle(2,3);\draw[dim](0,-.5)--(2,-.5)node[midway,below]{$b$};\draw[dim](-.5,0)--(-.5,3)node[midway,left]{$h$};\node at(1,3.6){闸门正视};
\end{scope}
\begin{scope}[xshift=5cm]
\fill[blue!10](-1,0)rectangle(0,4);\draw[guide](-1.4,4)--(2,4)node[right]{液面};
\draw[world,line width=1.5pt](0,0)--(0,3);\draw[dim](1,3)--(1,4)node[midway,right]{$z_0$};
\draw[dim](1,0)--(1,3)node[midway,right]{$h$};\node[left]at(-1,1.6){$\rho$};\node at(2.2,2){大气};
\end{scope}
\end{tikzpicture}}
\end{question}
\answerlink{s5:drill:2}
\begin{question}\begin{drill}[缩尺模型的摄影速率]\label{s5:drill:3}
自由下落的模型线度是实物的$1/49$，重力加速度相同。播放速率为每秒$24$帧，要让放大后的下落运动逼真，求拍摄帧率和对应速度比。
\end{drill}
\fig{fig:s5:drill:3:s}{题设：缩尺模型的摄影速率}{\begin{tikzpicture}[scale=.88]\draw[dim](0,0)--(0,3)node[midway,left]{$h$};\draw[dim](3,0)--(3,1.5)node[midway,left]{$h_m$};\fill(1,3)circle(3pt);\fill(4,1.5)circle(3pt);\draw[force](1,2.5)--(1,1.5)node[right]{$g$};\draw[force](4,1)--(4,0)node[right]{$g$};\node at(2,-.6){两图采用不同绘图比例，$h_m=h/49$};\end{tikzpicture}}
\end{question}
\answerlink{s5:drill:3}
\begin{question}\begin{drill}[车篷挡雨的速度三角形]\label{s5:drill:4}
车向右行驶，车篷后缘$A$正下方为$C$，后座边缘$B$在$C$左方，$AB$与竖直成$30^\circ$。车速$6\,\mathrm{km/h}$时雨相对车竖直，车速$18\,\mathrm{km/h}$时雨线沿$AB$。求雨对地速度。
\end{drill}
\fig{fig:s5:drill:4:s}{题设：车篷挡雨的速度三角形}{\begin{tikzpicture}[scale=.88]\draw[world](0,0)--(5,0);\draw[world](2,2)--(4.5,2);\draw[guide](2,2)--(2,0);\draw[guide](2,2)--(.845,0);\node[above]at(2,2){$A$};\node[below]at(2,0){$C$};\node[below]at(.845,0){$B$};\draw[world](2,1.4)arc(-90:-120:.6);\node at(1.55,1.05){$30^\circ$};\draw[flow](3.2,.6)--(4.5,.6)node[right]{$u$};\end{tikzpicture}}
\end{question}
\answerlink{s5:drill:4}
\begin{question}\begin{drill}[细线从圆柱展开的轨迹]\label{s5:drill:5}
半径$R$的固定圆柱上绕轻绳，总可展开长度$L$，一端固定于表面。质点从绳全绕紧处以速度$v$沿柱面外法线出发，绳始终绷紧。忽略重力，柱面碰撞完全弹性。求轨迹和一个完整往复周期。
\end{drill}
\fig{fig:s5:drill:5:s}{题设：细线从圆柱展开的轨迹}{\begin{tikzpicture}[scale=.88]\draw[world](1.7,0)circle(1.2);\fill(2.9,0)circle(3pt)node[below right]{$P_0$};\draw[flow](2.9,0)--(4.2,0)node[above]{$v$};\fill(.512,.1693)circle(2pt)node[left]{$A$};\node at(1.7,1.7){初态细线全绕紧};\node at(1.7,0){$O$};\draw[guide](1.7,0)--(1.7,-1.2)node[midway,right]{$R$};\node at(1.7,-1.7){从$P_0$到固定端$A$的绕绳长为$L$};\end{tikzpicture}}
\end{question}
\answerlink{s5:drill:5}
\begin{question}\begin{drill}[双秤砣称重的零点影响]\label{s5:drill:6}
杆秤用单秤砣读数准确。双秤砣称$8\,\mathrm{kg}$物体时读数$3\,\mathrm{kg}$，再用双秤砣称西瓜读数$6.5\,\mathrm{kg}$。求西瓜质量。
\end{drill}
\fig{fig:s5:drill:6:s}{双秤砣位于同一刻度处；零刻度与支点不必重合}{\begin{tikzpicture}

\draw[world](-1.5,0)--(4.8,0);\draw[world](-.25,-.3)--(0,0)--(.25,-.3)--cycle;\node[above]at(0,0){$O$};
\draw[world](-1.2,0)--(-1.2,-1);\draw[world](3.5,0)--(3.5,-1);\draw[world](3.1,-1)rectangle(3.9,-1.5);
\node at(3.5,-1.25){$2m$};\node[below]at(-1.2,-1){$M$};
\foreach \x in {1,1.5,2,2.5,3,3.5,4}{\draw(\x,0)--(\x,.15);}
\node[above]at(1,.15){零刻度};\node[above]at(3.5,.15){读数$q$};
\draw[dim](-1.2,.8)--(0,.8)node[midway,above]{$a$};\draw[dim](0,.8)--(1,.8)node[midway,above]{$c$};\draw[dim](1,.8)--(3.5,.8)node[midway,above]{$q\ell$};

\end{tikzpicture}}
\end{question}
\answerlink{s5:drill:6}
\begin{question}\begin{drill}[抽动小车时的圆柱接触]\label{s5:drill:7}
圆柱同时接触固定斜墙和水平小车，斜墙与水平成$\theta$且位于圆柱左上侧，车面向左抽动，速度$U<0$。圆柱重$G$、半径$R$、对轴转动惯量$I$，初始静止；墙、车面摩擦系数为$\mu_1,\mu_2$，在本模型中各处最大静摩擦系数与滑动摩擦系数取相同值。取$0<\theta<\pi/2$，讨论双接触时的运动约束和初始粘着条件。
\end{drill}
\fig{fig:s5:drill:7:s}{上方斜墙固定，下方车面左移；两个接触点的法向位置均固定}{\begin{tikzpicture}

\draw[world](-.5,.712)--(4.6,3.656);\draw[world](-.5,0)--(5,0);\draw[world](2,1)circle(1);
\draw[guide](-.5,.712)--(1,.712);\draw[world](.2,.712)arc(0:30:.7);\node at(.5,.96){$\theta$};
\fill(1.5,1.866)circle(2pt)node[above left]{$W,\mu_1$};\fill(2,0)circle(2pt)node[below]{$B,\mu_2$};
\draw[guide](2,1)--(3,1)node[midway,above]{$R$};\fill(2,1)circle(2pt)node[left]{$C$};
\draw[flow](4,-.7)--(2.8,-.7)node[left]{$U<0$};

\end{tikzpicture}}
\end{question}
\answerlink{s5:drill:7}
